\documentclass[11pt,a4paper]{article}

\usepackage[utf8]{inputenc}

\usepackage[T1]{fontenc}

\usepackage[polish]{babel}

\usepackage{lmodern}

\usepackage{geometry}

\geometry{margin=27mm}

\usepackage[hidelinks]{hyperref}

\title{Nasilenie objawów pourazowych a deklarowane korzystanie z generatywnej AI do wsparcia emocjonalnego: badanie przekrojowe}

\author{}

\date{Szkic naukowy, wersja 0.1.0}

\begin{document}

\maketitle

\begin{abstract}

Tematem pracy jest Nasilenie objawów pourazowych a deklarowane korzystanie z generatywnej AI do wsparcia emocjonalnego: badanie przekrojowe. Tekst systematyzuje uzasadnienie, hipotezy i sposób weryfikacji zagadnienia. Proponowany wskaźnik: Asocjacyjny OR z dwustronnym 95\% CI na 1 SD PCL-5 wobec binarnego Y\_E w zdefiniowanej populacji A\_X, zgodnie z SAP pracy. Nie przedstawia się niezweryfikowanych wyników ani fikcyjnych cytowań.

\end{abstract}

\section{Wprowadzenie}

Zachowania samoregulacyjne wymagają rozdzielenia stanu emocjonalnego, sytuacji zadaniowej i fizjologicznych korelatów. Zagadnienie Nasilenie objawów pourazowych a deklarowane korzystanie z generatywnej AI do wsparcia emocjonalnego: badanie przekrojowe wymaga odróżnienia badanego konstruktu od jego operacjonalizacji i kontekstowych źródeł zmienności. Trzeba porównać możliwe wyjaśnienia i ustalić granice interpretacji dowodów.

\section{Cel i pytania badawcze}

Celem pracy jest rozstrzygnięcie zagadnienia: Nasilenie objawów pourazowych a deklarowane korzystanie z generatywnej AI do wsparcia emocjonalnego: badanie przekrojowe. Proponowany wskaźnik: Asocjacyjny OR z dwustronnym 95\% CI na 1 SD PCL-5 wobec binarnego Y\_E w zdefiniowanej populacji A\_X, zgodnie z SAP pracy. Hipotezy i ich kierunki muszą zostać określone przed analizą danych lub krytyczną syntezą źródeł.

\section{Materiał i metody}

Zdefiniować warunki porównawcze, sposób oceny reakcji i czynniki zakłócające. Przed analizą należy określić jednostkę obserwacji, warunki kontrolne, kryteria wyłączenia oraz sposób ograniczania zmiennych zakłócających.  Przed rozpoczęciem badań należy udokumentować warunki etyczne, zgodę, prawa do danych i rzeczywistą dostępność narzędzi.

\section{Plan analizy}

Ocenić różnice między warunkami, zmienność indywidualną oraz wiarygodność sygnału. Kryteria wyboru modelu analitycznego, analiza niepewności i postępowanie z brakami danych wymagają uprzedniej specyfikacji.

\section{Wyniki}

W niniejszym pliku nie umieszczono wyników empirycznych, statystyk ani nieprzeprowadzonych analiz. Sekcja wymaga uzupełnienia po uzyskaniu i weryfikacji materiału dowodowego.

\section{Dyskusja i ograniczenia}

Ograniczenia obejmują możliwe błędy pomiaru, dobór próby, alternatywne mechanizmy i ograniczoną możliwość uogólniania. Korelacja nie dowodzi przyczynowości; poprawność techniczna nie jest dowodem znaczenia klinicznego.

\section{Wnioski}

Zdefiniowano ramy opracowania, ale nie sformułowano wniosków wykraczających poza dostępny stan dowodów.

\section*{Bibliografia}

Weryfikacja literatury i przypisanie cytowań do stwierdzeń są kolejnym etapem prac.

\end{document}

